\documentclass[11pt]{article}
\usepackage[T1]{fontenc}
\usepackage{textcomp}
\usepackage[margin=0.75in]{geometry}
\usepackage{amsmath,amssymb,amsthm}
\usepackage{array,booktabs}
\usepackage{hyperref}
\setlength{\parindent}{0pt}
\newtheorem{theorem}{Theorem}
\theoremstyle{definition}
\newtheorem{definition}{Definition}
\title{Flow Proof Artifact: Nat-plus-six-one}
\author{Generated by \texttt{flow doc proof}}
\date{Flow Proof Book}
\begin{document}
\maketitle
Six plus one is seven for natural numbers.\\[0.5em]
\textbf{Source.} peano — https://en.wikipedia.org/wiki/Peano\_axioms\\[0.5em]
\bigskip
\noindent{\large \textbf{Derived fact 1.} \textit{6 + 1 = 7} \hfill the natural numbers \cdot addition \cdot six plus one is seven}\par
\smallskip\noindent\textit{Coordinate: the natural numbers \cdot addition \cdot six plus one is seven}\par
\smallskip\noindent\textit{Source: peano}\par
\smallskip\noindent\textit{Built on: successor on the right via commutativity, for addition on the natural numbers, adding zero on the right does not change the number}\par
\medskip
\noindent\textbf{Goal.} 6 + 1 = 7\par
\noindent$\displaystyle six + \mathrm{succ}(0) = \mathrm{succ}(six)$\par
\medskip
\noindent
\begin{tabular}{@{} >{\raggedright\arraybackslash}p{0.06\textwidth} >{\raggedright\arraybackslash}p{0.47\textwidth} >{\raggedleft\arraybackslash}p{0.41\textwidth} @{}}
\textbf{\#} & \textbf{Proof} & \textbf{Mathematics} \\
\midrule
\textbf{1.} & We prove that six plus one is seven for addition on the natural numbers. & \\[0.45em]
\textbf{2.} & Let six = succ(succ(succ(succ(succ(succ(0)))))). & \\[0.45em]
\textbf{3.} & We invoke the derived fact governing addition on the natural numbers: successor on the right via commutativity, for addition on the natural numbers (instantiated for six, 0). & \\[0.45em]
\textbf{4.} & We invoke the derived fact governing addition on the natural numbers: adding zero on the right does not change the number (instantiated for six). & $\displaystyle six + 0 = six$ \\[0.45em]
\textbf{5.} & From step 2, step 3, and step 4, this implies six plus the successor of 0 equals the successor of six. Hence proven. & $\displaystyle six + \mathrm{succ}(0) = \mathrm{succ}(six)$ \\[0.45em]
\bottomrule
\end{tabular}
\label{thm:Nat-addition-six_plus_one_is_seven}
\par\medskip\hrule
\end{document}
